%%%%%%%%%%%%%%%%%%%%%%%%%%%%%%%%%%%%%%%%%
% Medium Length Professional CV
% LaTeX Template
% Version 2.0 (8/5/13)
%
% This template has been downloaded from:
% http://www.LaTeXTemplates.com
%
% Original author:
% Trey Hunner (http://www.treyhunner.com/)
%
% Important note:
% This template requires the resume.cls file to be in the same directory as the
% .tex file. The resume.cls file provides the resume style used for structuring the
% document.
%
%%%%%%%%%%%%%%%%%%%%%%%%%%%%%%%%%%%%%%%%%

%----------------------------------------------------------------------------------------
%	PACKAGES AND OTHER DOCUMENT CONFIGURATIONS
%----------------------------------------------------------------------------------------

\documentclass{resume} % Use the custom resume.cls style

\usepackage[left=0.75in,top=0.6in,right=0.75in,bottom=0.6in]{geometry} % Document margins
%\usepackage[colorlinks,linkcolor=red,anchorcolor=blue,citecolor=green]{hyperref}
\usepackage[colorlinks]{hyperref}
%\renewcommand{\baselinestretch}{2} \normalsize
%\usepackage{times}
\name{Tao Lin} % Your name

\address{State Key Lab of CAD\&CG, Zhejiang University \\ 310058 \\ Hangzhou, China} % Your address
%\address{123 Pleasant Lane \\ City, State 12345} % Your secondary addess (optional)
\address{http://nblintao.github.io/ \\ (86) 15957445452 \\ nblintao@gmail.com} % Your phone number and email

\begin{document}

\pagestyle{fancy}
\renewcommand{\headrulewidth}{0pt} 
\renewcommand{\footrulewidth}{0pt}


%----------------------------------------------------------------------------------------
%	EDUCATION SECTION
%----------------------------------------------------------------------------------------

\begin{rSection}{Education}

%{\bf Zhejiang University} \hfill { October 2012 - Present} \\ 
%B.Eng. in Computer Science \& Technology 
%\hfill {\em Hangzhou, China}
%\smallskip
%\\
\begin{rAwards}
	{Zhejiang University}
	{Oct. 2012 - Present}
	{B.Eng. in Computer Science and Technology} 
	{Hangzhou, China}
\end{rAwards}
{\contentSize
\begin{tabular}{ll}
	Overall GPA: & 3.96 / 4.00 (92.45 / 100)\\   %
	The third year GPA: & 4.00 / 4.00 (94.43 / 100) \\ %
%	Rank: & $\mathbf{1^{st}}$ in Qizhen Honors Class in Chu Kochen Honors College\\
%	 & Top 2\% among 196 students
	%Rank: & $\mathbf{1^{st}}$ in Qizhen Honors Class in Chu Kochen Honors College (Top 2\%)
	Rank: & $1^{st}$ in Qizhen Honors Class, Chu Kochen Honors College (Top 2\%)
\end{tabular}
}

\begin{rAwards}
	{Hong Kong University of Science and Technology (HKUST)}
	{Sept. 2015 - Feb. 2016}
	{Exchange Student for Final Year Project} 
	{Hong Kong}
\end{rAwards}


\end{rSection}



%----------------------------------------------------------------------------------------
%	PUCLICATION SECTION
%----------------------------------------------------------------------------------------

\begin{rSection}{Publications}
	
	\begin{rPublication}
		{Mobility Viewer: A Eulerian Approach for Studying Urban Crowd Flow}
		{Yuxin Ma, \textbf{Tao Lin}, Zhendong Cao, Chen Li, Wei Chen}
		{IEEE Transactions on Intelligent Transportation Systems. Accepted, 2015.}
		{\href{http://nblintao.github.io/pdf/Mobility\%20Viewer.pdf}{[PDF]}
			\href{https://youtu.be/m0HgyEsh0tI}{[Video]}}
		%http://www.cad.zju.edu.cn/home/vagblog/VAG_Work/tcmime.4202.4455.4674.pdf
		%\item A suite of interactive visualization techniques for exploring crowd flow volumes, directional dynamics and spatiotemporal patterns.
		%\item A visual analysis approach for inspecting human social relations and their physical movements.
	\end{rPublication}
	
	%------------------------------------------------
	
	\begin{rPublication}
		{TieVis: Visual Analytics of Evolution of Interpersonal Ties}
		{\textbf{Tao Lin}, Fangzhou Guo, Yingcai Wu, Wei Chen, Biao Zhu, Sidong Wang, Huamin Qu}
		%	{The 9th IEEE Pacific Visualization Symposium (PacificVis 2016). Submitted.}
		{IDAVis 2016: Workshop on Intelligent Data Analytics and Visualization. Submitted.}
		{	\href{http://nblintao.github.io/pdf/TieVis.pdf}{[PDF]}
			\href{https://youtu.be/xsr4YXOo1WA}{[Video]}
			\href{https://github.com/nblintao/TieVis}{[Code]}}
		
		%\item A new study of the evolution of interpersonal ties for a dynamic network and their co-evolution with the network structure and information diffusion.
		%\item An edge based analysis framework which helps users identify edges with similar trends, compare edges with different trends, and find the hidden patterns.
		%\item An interactive visualization system, which allows for multi-perspective exploration and analysis.
	\end{rPublication}
	
\end{rSection}

%----------------------------------------------------------------------------------------
%	WORK EXPERIENCE SECTION
%----------------------------------------------------------------------------------------

\begin{rSection}{Experience}

\begin{rExperience}
	{Visual Analytics Group, State Key Lab of CAD\&CG}
	{June 2013 - Aug. 2015}
	{Research Assistant}
	{Zhejiang University, China}
	{{\bf Research Advisor:} Professor Wei Chen {\scriptsize {\href{http://www.cad.zju.edu.cn/home/chenwei/}{[Homepage]}}}}
\item Administrated a computer cluster with 20 nodes for more than one year.
\item Processed large-scale data on the group's cluster with Hadoop and Spark for more than one year.
\item {\bf Project Mobility Viewer} - A visual analytics system supporting situation-aware understanding and visual reasoning of human mobility. 
{\scriptsize \href{http://nblintao.github.io/pdf/Mobility\%20Viewer.pdf}{[PDF]}%http://www.cad.zju.edu.cn/home/vagblog/VAG_Work/tcmime.4202.4455.4674.pdf
\href{https://youtu.be/m0HgyEsh0tI}{[Video]}}
\hfill {\em Core Team Member}

\begin{list}{$\cdot$}{\leftmargin=2em} % \cdot used for bullets, no indentation
	\itemsep -0.5em	\vspace{-0.5em} % Compress items in list together for aesthetics
\item Preprocessed tera-scale data on the cluster with Spark.
\item Designed a distributed algorithm to detect and eliminate ping-pong effects.
\item Designed and implemented the Flow Volume Analysis Subsystem.
\item Technique: MapReduce, Hadoop, Spark, Python, JavaScript (D3, WebGL), OpenSteetMap.

\end{list}


\item {\bf Project TieVis} - A visual analytics framework and system that enables interactive analysis and exploration of the dynamic changes of interpersonal ties.
	%	{May 2015 - Oct. 2015}{Project Leader}
{\scriptsize \href{http://nblintao.github.io/pdf/TieVis.pdf}{[PDF]}
\href{https://youtu.be/xsr4YXOo1WA}{[Video]}
\href{https://github.com/nblintao/TieVis}{[Code]}}
\hfill {\em Project Leader}
\begin{list}{$\cdot$}{\leftmargin=2em} % \cdot used for bullets, no indentation
	\itemsep -0.5em	\vspace{-0.5em} % Compress items in list together for aesthetics		
	\item Preprocessed the data with Python (Extract, Aggregate, PCA and MDS).	
	\item Designed and implemented the whole TieVis System.
	\item Technique: Python(NumPy, SciPy, scikit-learn), JavaScript.
\end{list}


\end{rExperience}

%------------------------------------------------

\begin{rExperience}
	{Multimedia Technology Research Center}
	{Sept. 2015 - Feb. 2016}
	{Visiting Intern}
	{HKUST, Hong Kong}
	{{\bf Supervisor:} Professor S.-H. Gary Chan {\scriptsize {\href{http://www.cse.ust.hk/~gchan/}{[Homepage]}}}}
\item Worked with a team that has been developing Streamphony, a next-generation streaming cloud for large-scale high bit-rate applications over the global Internet; helped turn research results into a practical system for commercial trials.
\item Implemented the forward error correction (FEC) module with Reed-Solomon codes.
\item Improved the quality and robustness of the video stream.
\item Currently developing a data analysis system for the proxy records with HBase.
\item Technique: C++ (STL, Boost (asynchronous, memory pool and smart pointer)), lock-free programming, code review, unit testing (Google Test, Google Mock), Hadoop (HBase).

\end{rExperience}

\end{rSection}

%----------------------------------------------------------------------------------------
%	AWARDS SECTION
%----------------------------------------------------------------------------------------

\begin{rSection}{Honors and Awards}

\begin{rAwards}{Google Excellence Scholarship}{2015}
	{One of 58 bachelor and master's students in Mainland China}{Google Inc.}
\end{rAwards}

\begin{rAwards}{National Scholarship}{2014 - 2015}{Top 1\% students}{Ministry of Education, China}
\end{rAwards}

%\begin{rAwards}{First Prize of Excellent Undergraduate Scholarship}{2013 - 2014, 2014 - 2015}{Top 3\% students}{Zhejiang University}
\begin{rAwards}{First-Class Scholarship for Outstanding Students}{2013 - 2014, 2014 - 2015}{Top 3\% students}{Zhejiang University}
\end{rAwards}

\begin{rAwards}{Honorable Mention, Interdisciplinary Contest in Modeling}{2015}{}{}
\end{rAwards}

\begin{rAwards}{Second Prize, Zhejiang University Programming Contest}{2014}{}{}
\end{rAwards}

\begin{rAwards}{Third Prize, ACM-ICPC China Zhejiang Provincial Programming Contest}{2014}{}{}
\end{rAwards}

\end{rSection}
%----------------------------------------------------------------------------------------
%	PROJECT SECTION
%----------------------------------------------------------------------------------------

\begin{rSection}{Selected Projects}



\begin{rSubsectionTag}
	{MiniSQL Database Management System}
	{Sept. 2014 - Nov. 2014}
	{Team Leader}
	{}
	{Course Project}
\item Designed and implemented a DBMS, which can create/delete/query tables/indexes/records by interpreting SQL.
\item Modules: Interpreter, API, Catalog Manager, Record Manager, Index Manager, Buffer Manager.
\item Technique: C++, Parser, Cache Replacement (LRU), B+ Tree.
\end{rSubsectionTag}

\begin{rSubsectionTag}
	{Visual Analysis of Library Data}
	{Apr. 2014 - May 2015}
	{Team Leader}
	{\href{https://youtu.be/MhnXvG_KWZs}{[Video]}
		\href{https://github.com/nblintao/Visual-Analysis-of-Library-Data}{[Code]}}
	{Student Research Training Program}
\item Promoted library users' decision-making by interacting with real records.  
\item Designed and implemented a visual analysis system.
\item Technique: JavaScript (D3, Bootstrap, ECharts and jQuery), Python (Django), Oracle Database.
\end{rSubsectionTag}


\begin{rSubsectionTag}
	{Precomputed Shadow Fields for Dynamic Scenes}
	{Mar. 2015 - July 2015}
	{Independent}
	{\href{https://github.com/nblintao/Precomputed-Shadow-Fields-for-Dynamic-Scenes}{[Code]}}
	{Course Project}
\item Implemented a system for rendering dynamic soft shadows based on the methodology outlined in a paper from SIGGRAPH 2005.
\item Course: Advances in Computer Graphics, Score: 94/100, Instructor: Kun Zhou.
\item Technique: C++ (DirectX), HLSL, precomputed radius transfer, spherical harmonics.
\end{rSubsectionTag}

%\begin{rSubsectionTag}
%	{Teaching Service System}
%	{March 2015 - July 2015}
%	{Group Leader}
%	{\href{https://github.com/nblintao/Score-Management-Subsystem}{[Subsystem Code]}
%	\href{https://github.com/nblintao/jwb}{[Full Code]}
%		}
%	{Course Project}
%\item Final project for Software Engineering, Score: 94, Instructor: Yue Chen.
%\item It is a large laboratory project developed by 30 students, using Python and JavaScript, with detailed documents and tests. I am the {\bf leader} of Score Management Subsystem.
%\end{rSubsectionTag}


\begin{rSubsectionTag}{A 3-D Auto-Stereoscopic Front Projection System}{July 2014 - May 2015}{Team Member}{\href{https://youtu.be/ExrAqQ9Radw}{[Video]}}{Challenge Cup Competition}
\item Worked with a team of optical engineering students to achieve glasses-free 3D by using 60 projectors.
\item Derived the core mathematical formulas by light field reconstruction principles.
\item Implemented the prototype with OpenGL.
\end{rSubsectionTag}



\end{rSection}


%
%%----------------------------------------------------------------------------------------
%%	SERVICE SECTION
%%----------------------------------------------------------------------------------------
%
%\begin{rSection}{Service}
%\begin{rAwards}{Volunteer}{April 2015}{The 8th IEEE Pacific Visualization Symposium}{Hangzhou, China}
%\end{rAwards}
%\end{rSection}

%----------------------------------------------------------------------------------------
%	TECHNICAL STRENGTHS SECTION
%----------------------------------------------------------------------------------------

\begin{rSection}{Technical Strengths}

\begin{tabular}{ @{} >{\bfseries}l @{\hspace{6ex}} l }
Computer Languages & C/C++, JavaScript, Python, Java, MATLAB, Verilog\\
Libraries \& Frameworks 
& Visualization (D3, ECharts) \\
& C++ unit testing (Google Test, Google Mock)\\
& Computer graphics \& vision (OpenGL, DirectX, WebGL, OpenCV)\\
Tools & Data science (Hadoop, Spark, Hive, HBase, MySQL)\\
& Software engineering (Unix Command, Makefile, Git, Phabricator)
\end{tabular}

\end{rSection}

%----------------------------------------------------------------------------------------
%	EXAMPLE SECTION
%----------------------------------------------------------------------------------------

%\begin{rSection}{Section Name}

%Section content\ldots

%\end{rSection}

%----------------------------------------------------------------------------------------

\end{document}




